\section{Thermal labels and the spectral update}\label{sec:thermal}

This section constructs the channels that will be attached to the spatial
pieces.  It also supplies two ways to pay for a change of channel: a
deterministic increase in a metric ball, and an observable event that is
unlikely under a comparison law.  The choices of channels and of input
membership sets will always be made without inspecting a sampled label.

\subsection{Positive Fourier kernels and form order}

\begin{definition}[Thermal forms]
A thermal form on the unit circle is a finite sum
\[
 A(y)=\sum_s\lambda_s(1-\operatorname{Re}\chi_s(y)),
 \qquad \lambda_s\ge0,\quad s\in\mathbb Z.
\]
We write $A\preceq A'$ if $A'=A+D$ for a thermal form $D$.  Define
\begin{equation}\label{eq:thermal-definitions}
 \begin{gathered}
 q_A(u)=\int_{\mathbb T}e^{-A(y)}\overline{\chi_u(y)}\,dy,
 \qquad F(A)=-\log q_A(0),\qquad K_A=e^{F(A)-A},\\
 \Delta_A(u)=\left(\int_{\mathbb T}|1-\chi_u(y)|^2K_A(y)\,dy\right)^{1/2},
 \qquad B_A(r)=\{u\in\mathbb Z:\Delta_A(u)\le r\}.
 \end{gathered}
\end{equation}
Thus $K_A$ is a probability density.  All circle integrals in this section
use normalized measure of total mass one.
\end{definition}

The covariance argument below belongs to the Griffiths--Ginibre
positive-correlation method \cite{Ginibre1970}. We include the circle
calculation, together with the form-order consequences needed for the
label updates.

\begin{lemma}[Thermal positivity and monotonicity]\label{lemma:thermal-monotonicity}
For every thermal form, $q_A$ is a nonnegative, even probability mass
function on $\mathbb Z$.  The coefficients
$\int\chi_uK_A=q_A(u)/q_A(0)$ are real and nonnegative.  For all integers
$u,v$,
\[
 \operatorname{Cov}_{K_A}(\operatorname{Re}\chi_u,
                          \operatorname{Re}\chi_v)\ge0.
\]
Consequently these characteristic coefficients increase in form order,
$\Delta_A(u)$ decreases in form order, and
$\Delta_A(u+v)\le\Delta_A(u)+\Delta_A(v)$.  The function $F$ increases in
form order.  For every thermal form $D$, the increment
$F(A+D)-F(A)$ decreases when $A$ increases in form order.

If $P=\sum_jP_j$ is a finite sum of thermal forms and
$P_{\rm lo}\preceq P\preceq P_{\rm hi}$, then, for
arbitrary shifts $z_j$ and every integer $u$,
\begin{equation}\label{eq:thermal-envelope}
 \left|\int\chi_u(y)\prod_j e^{-P_j(z_j-y)}\,dy\right|
 \le q_P(u)
 \le e^{-F(P_{\rm lo})}\int\chi_u K_{P_{\rm hi}}.
\end{equation}
\end{lemma}

\begin{proof}
Expand each factor as
\[
 e^{-\lambda(1-\operatorname{Re}\chi_s)}
 =e^{-\lambda}
   \exp\bigl((\lambda/2)\chi_s\bigr)
   \exp\bigl((\lambda/2)\chi_{-s}\bigr).
\]
Its Fourier coefficients are nonnegative, with sum one.  The expansions
are absolutely summable, as is their finite product.  Evaluation at zero
therefore gives $\sum_uq_A(u)=1$; evenness gives $q_A(-u)=q_A(u)$.

For completeness, write $c_s(y)=\cos(2\pi sy)$ and
$Z_A=\int\exp(\sum_s\lambda_sc_s)$.  Express the covariance using two
independent copies as half the expectation of the product of their
differences.  The map $(x,z)\mapsto(x+z,x-z)$ pushes Haar probability on
the two-dimensional torus to Haar probability.  The two differences are
$-2\sin(2\pi ux)\sin(2\pi uz)$ and the analogous expression for $v$,
whereas the unnormalized product density is
$\exp(2\sum_s\lambda_sc_s(x)c_s(z))$.  Expanding this last exponential
gives the identity
\[
 \operatorname{Cov}_{K_A}(c_u,c_v)
 =\frac{2}{Z_A^2}
 \sum_{(n_s)\ge0}\left(\prod_s\frac{(2\lambda_s)^{n_s}}{n_s!}\right)
 \left(\int\sin(2\pi ux)\sin(2\pi vx)
                 \prod_sc_s(x)^{n_s}\,dx\right)^2.
\]
Absolute convergence justifies the expansion and proves nonnegativity.
Differentiation under the integral yields
\[
 \frac{\partial}{\partial\lambda_v}\int c_uK_A
 =\operatorname{Cov}_{K_A}(c_u,c_v),\qquad
 \frac{\partial F}{\partial\lambda_u}=\int(1-c_u)K_A.
\]
The latter derivatives are nonnegative and decrease in every nonnegative
coefficient direction.  Integrating these derivative statements proves
the assertions about $F$ and its increments.  Since
$\Delta_A(u)^2=2(1-\int c_uK_A)$, the metric monotonicity follows as well.
Subadditivity follows from
$1-\chi_{u+v}=(1-\chi_u)+\chi_u(1-\chi_v)$ and the triangle inequality
in $L^2(K_A)$.

Translation multiplies each Fourier coefficient of $e^{-P_j}$ by a
unit complex number.  Absolute values in the convolution of these
coefficients are therefore bounded by the corresponding convolution of
the centered coefficients, which is $q_P(u)$.  Finally
$q_P(u)=e^{-F(P)}\int\chi_uK_P$, and both required comparisons follow
from the monotonicities already proved.
\end{proof}

For a real interval $p$ of length at most one, set
\[
 R_p=|p|^{-1},\qquad
 S_p(y)=R_p^2(1-\operatorname{Re}\chi_1(y)).
\]
The interval may be viewed on the circle when integrating.  We will use
the following bounds, also with any larger spatial scaling in place of
$R_p$:
\begin{equation}\label{eq:thermal-spatial}
 \begin{gathered}
 F(S_p)=\log R_p+O(1),\qquad
 \int(1-\operatorname{Re}\chi_1)K_{A+S_p}\le C R_p^{-2},\\
 q_{S_p}(u)\ge cR_p^{-1}\quad (|u|\le cR_p),\qquad
 \Delta_{A+S_p}(u)\le C|u|/R_p.
 \end{gathered}
\end{equation}
Indeed, if $\|y\|$ denotes distance to the nearest integer, then
$8\|y\|^2\le1-\cos(2\pi y)\le2\pi^2\|y\|^2$.  Gaussian comparison
gives $\int e^{-S_p}\asymp R_p^{-1}$ and
$\int(1-\operatorname{Re}\chi_1)K_{S_p}\le CR_p^{-2}$.  Monotonicity
gives the second bound for $A+S_p$.  The elementary inequality
$|1-\chi_u|\le|u||1-\chi_1|$ gives the last bound and, for sufficiently
small $|u|/R_p$, gives $\int\chi_uK_{S_p}\ge1/2$, proving the remaining
lower bound.  In particular $\int S_pK_{A+S_p}\le C$.

\subsection{Channels, scales, and the cost of tightening a form}

Take $b_0$ dyadic with $h^{-16}/2<b_0\le h^{-16}$.  The computation
levels are the dyadic $b$ from $b_0/8$ to $1$, and
\begin{equation}\label{eq:thermal-parameters}
 k_b=A_0\frac d h(1+\log(1/b))^2,\qquad
 r_{\rm sep}=h^{-30},\qquad \sigma=h^{-40}.
\end{equation}
The constants will be fixed in the order stated below.  Every estimate
holds for sufficiently large $d$, uniformly in the spatial depth and in
any frequency cutoff.

A label of type $A$ is $Z=Y+\theta$ modulo one, where the error $\theta$
has density $K_A$.  Initially $A=0$, so a dummy independent uniform label
may be drawn.  Only the current main channel is used in the spectra and
transforms.  A successful update adds a spatial form and a metric form,
$A^+=A+S_p+D_i$.  The errors used for different label draws are
independent given $Y$.  Their parameters are selected from spatial
prefixes and earlier input membership events.  In particular neither a
sampled old label nor a sampled fresh label selects a membership set,
channel, or test.  Conditional on $Y$ and its discrete parameter path,
the errors thus retain their prescribed independent laws.

We will need to tighten the forms of past labels at small free-energy
cost and to keep their own-error energies $\int A K_A$ small. Both estimates
follow by choosing each metric addition on a sufficiently flat interval of
free energy. We require, with $A_{\rm base}$ the form immediately before
the addition and including its new spatial term,
\begin{equation}\label{eq:source-5}
 F(A_{\rm base}+h^{140}D_i)-F(A_{\rm base}+D_i/2)
 \le d h^{-400}.
\end{equation}
Before a requested update assume that fewer than $C_0h$ main additions
have occurred and that
\begin{equation}\label{eq:source-6}
 F(A)\le\log R_*+C_*d.
\end{equation}
Here $p_*$ is the interval of the most recent addition, $R_*=R_{p_*}$,
and the current $p$ is contained in $p_*$.  Before the first addition
$p_*$ is the top interval, whose length lies in $[1/2,1]$.  The procedure
in \Cref{sec:procedure} will establish these hypotheses inductively.

Fix a sufficiently large numerical $C_s$ and put
$\widetilde A=h^{130}A+C_sS_p$.  Then
\begin{equation}\label{eq:source-7}
 \begin{split}
 F(\widetilde A)-F(A+S_p)
  &\le C+CC_0h\bigl(dh^{-400}+C\log h\bigr),\\
 \int A K_A&\le CC_0h\bigl(dh^{-400}+1\bigr).
 \end{split}
\end{equation}
To prove the first line, enlarge the constituent forms one at a time.
For an old metric term $D_i$, decreasing increments bound its cost of
tightening by the left side of \Cref{eq:source-5}: the other terms already
contain its original $A_{\rm base}$.  Tightening an old spatial term
costs at most
$F(h^{130}S_{p_i})-F(S_{p_i})\le C(1+\log h)$ by
\Cref{eq:thermal-spatial}.  Tightening the current newly inserted $S_p$
to $C_sS_p$ costs $O(1)$.  For the energy bound, concavity along the
$D_i$ direction and \Cref{eq:source-5} give
\[
 \int D_iK_A\le
 2\{F(A_{{\rm base},i}+D_i)-F(A_{{\rm base},i}+D_i/2)\}
 \le2dh^{-400}.
\]
Each spatial constituent contributes $O(1)$ by
\Cref{eq:thermal-spatial} and monotonicity.  This energy argument uses
the count bound and \Cref{eq:source-5}, but does not use
\Cref{eq:source-6}; it remains valid, with the same type of bound, just
after a final allowed addition that breaches a later stopping cap.

The spatial increment itself obeys
\begin{equation}\label{eq:thermal-spatial-increment}
 F(A+S_p)\le F(A)+\log(R_p/R_*)+C.
\end{equation}
After an addition, $A$ contains $S_{p_*}$, so
$q_A(u)\ge ce^{-F(A)}$ for $|u|\le cR_*$.  Convolution with $q_{S_p}$
over those integers gives
$q_{A+S_p}(0)\ge ce^{-F(A)}R_*/R_p$.
Initially use the term $u=0$ alone and $1\le R_*\le2$.

\subsection{From a separated spectrum to a finite compression problem}

Consider a subdensity $a$ at an interval $p$, with
\[
 0\le a\le e^d,\qquad e^{-d}\le m=\mathbb E_pa\le2\alpha,
 \qquad \alpha>0.
\]
Here $\mathbb E_p$ is uniform spatial average, and $\mathcal P$ is the
law of $Y$ with density $a/m$ relative to that average.  For either
prospective child $p'$ define
\begin{equation}\label{eq:thermal-child-moments}
 \begin{split}
 M_{p'}(z)&=\mathbb E_{p'}[aK_A(z-Y)]/\alpha,\\
 f_{p'}^u(z)&=\mathbb E_{p'}[a\chi_u(Y)K_A(z-Y)]/\alpha,\\
 G_{p'}^u&=\int |f_{p'}^u(z)|^2/M_{p'}(z)\,dz.
 \end{split}
\end{equation}
The integrand is zero at zero mass.  An integer is high at $b$ when
some $G_{p'}^u\ge b^2$.  A test requests an update if there is a
collection of $\lceil e^{k_b}\rceil$ high integers whose pairwise
differences avoid
\begin{equation}\label{eq:thermal-separation}
 B_A(r_{\rm sep})+\{j\in\mathbb Z:|j|\le\sigma R_p\}.
\end{equation}
One offending level and one collection are used at a time.

Keep at least half this collection having a common high child, and snap
each retained integer to a nearest multiple $gx$, where
\[
 g=\max(1,\lfloor\sigma R_p/20\rfloor).
\]
The snapping error is at most $\sigma R_p/20$, and is zero when $g=1$.
The snapped indices are distinct, since otherwise the original
difference would violate \Cref{eq:thermal-separation}.  Denote by $S$ the resulting set of integer indices $x$.  For each $x\in S$ choose a measurable complex function
$v_x$, with $|v_x|\le1$, so that
\begin{equation}\label{eq:source-8}
 \operatorname{Re}\mathbb E_{\mathcal P}
  [\mathbf1_{\rm child}v_x(Z)\chi_{gx}(Y)]\ge\beta,
 \qquad \beta=b^2/100.
\end{equation}
The expectation here includes the type-$A$ channel.  To see the bias,
$|f_{p'}^u|\le M_{p'}$ implies $\int|f_{p'}^u|\ge G_{p'}^u$.
Choosing the polarizing sign gives bias
$(\alpha/(2m))\int|f_{p'}^u|\ge b^2/4$ before snapping.  A constant
phase rotation at the center of $p$ makes the error of replacing $u$
by $gx$ at most $C\sigma$.  This is smaller than the required slack
because $\sigma=o(b^2)$ uniformly over the computation levels.

We apply \Cref{lemma:compression} with
\[
 d\omega=K_{\widetilde A}\,dy,\qquad
 q(x)=\frac{q_{\widetilde A}(gx)}{N_g},\qquad
 N_g=\sum_{x\in\mathbb Z}q_{\widetilde A}(gx).
\]
The spatial frequency-one term makes every Fourier coefficient strictly
positive: in its exponential expansion, every integer frequency occurs
with positive coefficient.  Therefore $q$ is strictly positive.  The
grid identity and Gaussian comparison give
\begin{equation}\label{eq:thermal-grid-normalizer}
 N_g=\frac1g\sum_{l=0}^{g-1}e^{-\widetilde A(l/g)},\qquad
 \frac1g\le N_g\le\frac Cg.
\end{equation}
Indeed $g\le R_p$, and the summands are bounded by
$\exp(-cR_p^2\|l/g\|^2)$.  Moreover $R_p/g\le C/\sigma$, also when
$g=1$ by its definition.  Thus \Cref{eq:source-6,eq:source-7,eq:thermal-spatial-increment}
imply
\[
 F_0=-\log q(0)=F(\widetilde A)-\log g+O(1)\le D_0d
\]
for a fixed sufficiently large $D_0$ depending only on the fixed caps.
Evenness gives precisely
$q(x)/q(0)=\int\chi_{gx}\,d\omega$.

We verify both remaining uniform hypotheses of compression.

\paragraph{Sumset size.}
We use the dissociated-spectrum method underlying Chang's lemma
\cite[Section~3]{Chang2002}, in its entropy and positive-product form
\cite[Section~3.2]{Lee2017}. The following calculation includes the side
information carried by the label: the comparison law makes that
information independent of the uniform circle coordinate.

The law of $gY$ modulo one has density at most
$e^{2d}(1+R_p/g)$: its preimage count in $p$ is at most $g/R_p+1$,
and the density of $Y$ relative to circle length is at most $R_pe^{2d}$.
Also
\begin{equation}\label{eq:thermal-label-information}
 \mathrm I(Y;Z)\le F(A)-\log R_*+C.
\end{equation}
Compare each $K_A(z-Y)$ to the fixed density
$K_{S_{p_*}}(z-\operatorname{center}(p_*))$.  The average relative
entropy is
\[
 F(A)-F(S_{p_*})-\int AK_A
 +\mathbb E S_{p_*}(\theta+Y-\operatorname{center}(p_*)).
\]
The last expectation is $O(1)$: use
$S(v+w)\le2S(v)+2S(w)$, the spatial energy bound, and
$|Y-\operatorname{center}(p_*)|\le|p_*|/2$.  In the initial case $A=0$
and $R_*\le2$, the same bound is immediate.  Dropping the nonpositive
term proves \Cref{eq:thermal-label-information}.

Let $Q$ denote the side information consisting of $Z$ and the child
bit.  Data processing and its one-bit entropy imply
\[
 D\bigl((gY,Q)\,\|\,\text{uniform}\otimes\mathcal L(Q)\bigr)
 =D(gY\|\text{uniform})+\mathrm I(gY;Q)\le C(C_*)d.
\]
If $\Lambda\subset S$ is dissociated, meaning that no nonzero
coefficient choice from $\{0,\pm1\}$ sums to zero, form the positive
product
\[
 \prod_{x\in\Lambda}
 \left(1+c\beta\operatorname{Re}
       [\mathbf1_{\rm child}v_x(Z)\chi_x(gY)]\right).
\]
For sufficiently small numerical $c$, the factors lie in $[1/2,3/2]$.
Under the reference law the first coordinate is uniform and independent
of $Q$; expansion and dissociation make the expectation of the product
exactly one.  Relative entropy is at least the true expected logarithm
of this product.  The inequality $\log(1+s)\ge s-Cs^2$, together with
\Cref{eq:source-8}, bounds the expected logarithm of each factor below
by $c'\beta^2$.  Hence $|\Lambda|\le C(C_*)d/\beta^2$.
A maximal dissociated subset spans $S$ by coefficients $0,\pm1$:
otherwise adjoining an unspanned element preserves dissociation.
Consequently
\[
 |nS|\le(2n+1)^{|\Lambda|},\qquad
 \log|nS|\le D_0d h^{P_0}\log d\quad(1\le n\le d^2)
\]
for fixed $D_0,P_0$, because $\beta\ge c h^{-32}$.

\paragraph{The uniform moment cap.}
For any probability $\nu$ on $S$ set
\begin{equation}\label{eq:thermal-moment-cost}
 \begin{split}
 E_T(\nu)&=-\log\mathbb E_\nu
 H_q\left(\sum_{i=1}^TX_i-\sum_{i=1}^TX_i'\right),\\
 e^{-E_T(\nu)}&=
 \int\left|\sum_x\nu(x)\chi_{gx}(y)\right|^{2T}K_{\widetilde A}(y)\,dy,
 \end{split}
\end{equation}
where the $2T$ letters are independent with common law $\nu$.  Then
\begin{equation}\label{eq:thermal-moment-cap}
 0\le E_T(\nu)\le D_0d\qquad\hbox{for every such }\nu.
\end{equation}
Here is a direct proof of the upper bound.  Restrict $Z$ to the domain
$\Omega$ at circle distance at most $h^{40}/R_*$ from $p$.  The channel
spatial variance bounds the discarded probability by $Ch^{-80}$
(the assertion is trivial if $\Omega$ is the full circle).
This loses $o(\beta)$ in \Cref{eq:source-8}.  Summing the biases with
weights $\nu$ and applying Jensen on a subprobability measure gives
\[
 \mathbb E_{\mathcal P}
 \left[\mathbf1_{\rm child}\mathbf1_{\{Z\in\Omega\}}
 \left|\sum_x\nu(x)v_x(Z)\chi_{gx}(Y)\right|^{2T}\right]
 \ge(\beta/2)^{2T}.
\]
For an upper bound discard the child indicator.  At fixed $z$, first
majorize the nonnegative moment integrand using the spatial density
bound $R_pe^{2d}$ and
\[
 \mathbf1_p(y)\le
 \exp\bigl(C-S_p(y-\operatorname{center}(p))\bigr).
\]
Only after this majorization expand the even power and the shifted
thermal weights.  By \Cref{eq:thermal-envelope}, the Fourier
coefficients of those weights are bounded absolutely by
\[
 Cq_{A+S_p}(u)\le
 Ce^{-F(A+S_p)}\int\chi_u K_{\widetilde A}.
\]
The coefficients $v_x(z)$ have modulus at most one.  The absolute
coefficient sum in the even-power expansion therefore averages this
envelope at $u=g(\sum X_i-\sum X_i')$ over the independent letters in
\Cref{eq:thermal-moment-cost}.  Using $|\Omega|\le Ch^{40}/R_*$, the
joint moment is at most
\[
 C R_p e^{2d+F(A)}\frac{h^{40}}{R_*}
       e^{-F(A+S_p)}e^{-E_T(\nu)}.
\]
Now $F(A+S_p)\ge\log R_p-C$,
$F(A)\le\log R_*+C_*d$, and
$T\log(2/\beta)=o(d)$.  Comparison with the lower moment proves
\Cref{eq:thermal-moment-cap} after increasing the fixed $D_0$.
This argument uses a density bound only in a nonnegative integral; it
does not assert a Fourier coefficient bound for $a$.

\subsection{The iterative update and the deterministic alternative}

\begin{lemma}[Spectral update]\label{lemma:spectral-update}
Fix $C_0,C_*\ge1$.  Let the current form $A$ arise from the initial
form zero by fewer than $C_0h$ additions of spatial and metric forms,
each metric satisfying \Cref{eq:source-5}.  Let $p\subset p_*$, where
$p_*$ is the latest addition interval, or the top interval of length in
$[1/2,1]$ if there has been no addition, and assume
$F(A)\le\log R_{p_*}+C_*d$.  On $p$ suppose
$0\le a\le e^d$ and $e^{-d}\le m=\mathbb E_pa\le2\alpha$, with
$\alpha>0$.  Suppose there is a request at a computation level $b$ in
the sense of \Cref{eq:thermal-separation}, using the prospective child
moments in \Cref{eq:thermal-child-moments} and the type-$A$ channel.
There are numerical constants $C_1,C_2$
and a sufficiently small numerical $\varepsilon>0$ with the following
property.  After the caps are fixed, choose fixed sufficiently large
$D_0,P_0$, then $A_0$, and finally take $d$ sufficiently large.
A finite deterministic construction supplies a parameter
\[
 d/h^2\le t\le D_0d,\qquad t\ge\varepsilon k_b,
\]
a thermal metric form $D_i=\lambda B$, and an input-only success set.
On success set $A^+=A+S_p+D_i$.  The new metric satisfies
\Cref{eq:source-5} and
\begin{equation}\label{eq:source-9}
 \begin{gathered}
 A^+=A+S_p+D_i,\qquad \int B K_{A^+}\le C_1t/\lambda,\\
 \lambda\ge dh^{200},\qquad h^{140}\lambda\le L,
 \end{gathered}
\end{equation}
and $F(A^+)-F(A+S_p)\le C_1t$.
There are two alternatives:
\begin{enumerate}[label=\textup{(\roman*)}]
 \item A regular test has success probability $\delta\ge\beta/2$ under
 the current input law, where $\beta=b^2/100$.  It also supplies a
 probability $\nu$ on a set $S'$ and a center $x_0\in S$ for which
 $E_T(\nu)\ge t/4$.
 \item A cheap update succeeds deterministically, has
 $t=\varepsilon k_b$, and uses the uniform probability on $S'$.
\end{enumerate}
In both alternatives $\log|S'|\ge0.8k_b$ and
$B=\sum_x\nu(x)(1-\operatorname{Re}\chi_{g(x-x_0)})$.
Track a radius starting at $r_{\rm sep}/4$ and increasing by
$\eta=r_{\rm sep}/(10C_0h)$ at every main addition.  As long as the
count cap applies, the gap $F(A)-\log|B_A(r)|$ starts at zero, stays
bounded below by a numerical negative constant, increases by at most
$C_2t$ at a regular success, and decreases by at least $k_b/2$ at a
cheap success.  On failure there is no main addition and no label draw.
All these conclusions are independent of spatial depth and frequency
cutoff.
\end{lemma}

\begin{proof}
\textbf{Selecting the subset and its weight.}
On $S$ and then on its successive restrictions set
\[
 t=\max\left(\varepsilon k_b,
                 \max_{\nu\text{ supported on the current set}}E_T(\nu)
       \right).
\]
The maximum exists: the simplex is compact, its moment integral is
continuous and strictly positive, and therefore its negative logarithm
is continuous.  For large $d$, the floor is between $d/h^2$ and $D_0d$;
the upper bound follows also for the maximum from
\Cref{eq:thermal-moment-cap}.  Apply \Cref{lemma:compression} with
$E=t$ to obtain a subset $S'$ and center $x_0$. Its conclusion holds
for every probability on this fixed $S'$, so the weight may now be chosen
by the following test.
If the maximum of $E_T$ on $S'$ is at least $t/4$, keep a maximizing
witness and terminate with a regular test.  If that maximum is smaller
than $t/4$ and $t=\varepsilon k_b$, take uniform $\nu$ on $S'$ and
terminate with a cheap update.  Otherwise replace the current set by
$S'$ and repeat.  Every nonterminal iteration reduces $t$ by a factor
of at least four, except possibly for the last transition to the fixed
floor.  Thus the process terminates in finitely many iterations, and
the sum of the logarithmic cardinality losses is at most
\[
 C(D_0)\frac{d^2}{\varepsilon k_bh^2}.
\]
The initial child restriction lost at most $\log2$ in logarithmic
cardinality.  Since $k_b\ge A_0d/h$, choosing $A_0$ sufficiently large
in terms of $D_0,P_0,\varepsilon$ makes the final loss less than
$0.2k_b$, proving $\log|S'|\ge0.8k_b$.

\textbf{Choosing a flat coefficient range.}
For the last center and chosen probability, the compression conclusion
is
\[
 F(\widetilde A+LB)-F(\widetilde A)\le Ct.
\]
Put $A_{\rm base}=A+S_p$. Monotonicity of $F$ and
$A_{\rm base}\preceq\widetilde A$ give
\[
 F(A_{\rm base}+LB)-F(A_{\rm base})
 \le \bigl[F(\widetilde A+LB)-F(\widetilde A)\bigr]
       +\bigl[F(\widetilde A)-F(A_{\rm base})\bigr].
\]
The first bracket is controlled by compression and the second by the
tightening bound in \Cref{eq:source-7}. Hence
\begin{equation}\label{eq:thermal-volume-cost}
 F(A_{\rm base}+LB)-F(A_{\rm base})\le C_1t.
\end{equation}
The constant $C_1$ is numerical: the volume constant of compression is
numerical, and the cap-dependent error in \Cref{eq:source-7} is $o(t)$
after the fixed caps and $A_0$ have been selected.

To find a flat range, set
\[
 \lambda_j=dh^{200}(2h^{140})^j,
 \qquad
 N=\left\lfloor
 \frac{\log(2L/(dh^{200}))}{\log(2h^{140})}\right\rfloor.
\]
For $0\le j<N$, the intervals
$[\lambda_j/2,h^{140}\lambda_j]$ are adjacent and lie in
$[dh^{200}/2,L]$.  Since $\log(L/d)=w=e^{h/4}$,
$N\asymp e^{h/4}/\log h>C_1D_0h^{400}$ for sufficiently large $d$.
The sum of the free-energy increments on these intervals is bounded
by \Cref{eq:thermal-volume-cost}, hence by $C_1D_0d$.  At least one
increment is at most $dh^{-400}$.  Take its $\lambda$ and
$D_i=\lambda B$; this is \Cref{eq:source-5}.  Concavity along the
$B$ direction gives
\[
 \int BK_{A_{\rm base}+\lambda B}
 \le\frac{F(A_{\rm base}+\lambda B)-F(A_{\rm base})}{\lambda}
 \le C_1t/\lambda,
\]
which proves \Cref{eq:source-9} and the metric increment bound.

\textbf{The input-only success set.}
For a regular test define the success set by
\begin{equation}\label{eq:thermal-success-set}
 \left|\sum_x\nu(x)\chi_{gx}(Y)
                  \mathbb E[v_x(Z)\mid Y]\right|\ge\beta/2.
\end{equation}
The expectation uses the specified type-$A$ channel, rather than the
sampled old observation.  The expression has modulus at most one, and
\Cref{eq:source-8}, averaged with $\nu$, gives its expected modulus at
least $\beta$.  Thus its probability of exceeding $\beta/2$ is at
least $\beta/2$.  In the cheap alternative take the entire input as
success.  The data $a,p,\alpha,A$ determine all these sets and choices.

\textbf{The change in the gap potential.}
It remains to prove the gap assertions. Since $r\le r_{\rm sep}/2$
through the count cap, summability of $q_A$ gives
\[
 |B_A(r)|e^{-F(A)}(1-r^2/2)\le1,
 \qquad F(A)-\log|B_A(r)|\ge\log(1-r^2/2)\ge-C.
\]
The balls are finite by the same inequality.  Initially $A=0$ and
$B_0(r)=\{0\}$, so the gap is zero.
The spatial addition raises $F$, but also allows short integer translates
of the old ball inside a slightly larger new ball. We compare these changes
using the same overlap multiplicity. Write
\[
 U=B_A(r),\qquad
 D_p=\mathbb Z\cap[-\sigma R_p/10,\sigma R_p/10],\qquad
 l=|(U-U)\cap(D_p-D_p)|.
\]
Here $l\ge1$.  On this intersection $\Delta_A\le2r$ and
$|u|\le\sigma R_p/5$, so convolution and
\Cref{eq:thermal-spatial} imply
\begin{equation}\label{eq:thermal-packing-pair}
 F(A+S_p)-F(A)\le\log R_p-\log l+C,
 \qquad |U+D_p|\ge |U||D_p|/l.
\end{equation}
For the second inequality, each representation of a fixed element of
$U+D_p$ differs from one fixed representation by an element of the
intersection defining $l$, so its multiplicity is at most $l$.

After either addition, $U+D_p$ lies in $B_{A^+}(r+\eta/2)$, since
monotonicity preserves the old radius and
$\Delta_{A^+}(j)\le C|j|/R_p\le C\sigma=o(\eta)$ for $j\in D_p$.
On a cheap success, uniformity of $\nu$ and \Cref{eq:source-9} give
\[
 \frac1{|S'|}\sum_{x\in S'}\Delta_{A^+}(g(x-x_0))^2
 \le 2C_1t/\lambda.
\]
At least half the $x$ therefore have
$\Delta_{A^+}(g(x-x_0))\le\sqrt{4C_1t/\lambda}
\le C(D_0)h^{-100}=o(\eta)$.  Their translates
$g(x-x_0)+U+D_p$ lie in the new tracked ball $B_{A^+}(r+\eta)$.
They are disjoint.  A collision would put $g(x-x')$ in
$U-U+D_p-D_p$; after restoring the two snapping errors, the difference
of the original frequencies would lie in
\[
 B_A(2r)+\{j\in\mathbb Z:|j|\le3\sigma R_p/10\},
\]
contrary to \Cref{eq:thermal-separation}.

Since $R_p/|D_p|\le C/\sigma$, including the case $D_p=\{0\}$,
\Cref{eq:thermal-packing-pair} cancels the same $\log l$ in the
free-energy and cardinality costs.  On a regular success the gap
increase is at most $C_1t+O(\log h)\le C_2t$ for a numerical $C_2$.
On a cheap success it is at most
\[
 C_1\varepsilon k_b-\log(|S'|/2)+O(\log h)
 \le C_1\varepsilon k_b-0.8k_b+O(\log h)\le-k_b/2,
\]
on choosing $\varepsilon$ sufficiently small numerically and then
taking $d$ sufficiently large.  These choices precede the cap choices;
only the size-loss bound and the small-error thresholds require the
subsequent choices of $A_0$ and $d$.
\end{proof}

\subsection{Evidence under the comparison law and under truth}

At every success draw a new main label $V$ of type $A^+$.  On a regular
success also draw $N_h=\lceil h^{120}\rceil$ fresh, independent
type-$A$ clones $Z'$.  They serve only the evidence calculation.
These draws are made even when this very success will breach a cap and
transfer the piece out of live processing.  A failure causes neither a
new main draw nor an evidence draw.

The comparison experiment starts with $Y$ uniform on the top interval.
Spatial bits continue to observe $Y$, and all label draws use the same
channels prescribed on their recorded branches.  At a membership test,
however, success is an independent coin with probabilities
$\delta,1-\delta$ computed from the current true input law before
conditioning on sampled labels.  For a fixed recorded spatial and
membership branch, the computed state is fixed.  Its membership coin
probabilities therefore contribute no $y$-dependent likelihood.  Given
the full past observations at a test in $p$, the comparison posterior
has density on $p$ proportional to
\begin{equation}\label{eq:thermal-reference-posterior}
 \exp\left(-\sum_jA_j(z_j-y)\right).
\end{equation}
The sum includes past main labels and past evidence clones.  Each
$A_j\preceq A$, and their number is at most
$1+C_0h(1+N_h)$.  This statement conditions on the observed branch,
not on the actual truth-membership event within the reference law.

Set $s_d=dh^{-100}$.  A past is admissible when some
$y_0\in\overline p$ satisfies
\[
 \sum_jA_j(z_j-y_0)\le s_d.
\]
This is an observable condition: the continuous function in the display
attains its infimum on the compact interval closure.

\begin{proposition}[Regular evidence]\label{prop:thermal-evidence}
Consider any finite procedure starting with the type-zero label, making
at most $\lfloor C_0h\rfloor$ successful main additions and using the
updates of \Cref{lemma:spectral-update} with its hypotheses at every
attempt.  Require all discrete choices to depend only on spatial
prefixes and input memberships, with independent label errors given
$Y$ and its parameter path, and include all the success draws specified
above.  Then an admissible regular attempt
has the following conditional reference bound, given its past and its
success coin.  The observable event
\begin{equation}\label{eq:source-10}
 \begin{gathered}
 \text{for some fresh clone }Z',\qquad A(Z'-Z)\le4s_d,\\
 \left|\sum_x\nu(x)v_x(Z')\chi_{gx}(V)\right|\ge\beta/4.
 \end{gathered}
\end{equation}
has probability at most $e^{-t/8}$, where $Z$ is the old main label.
Under the true input law, all attempted pasts are admissible and
\Cref{eq:source-10} occurs at every regular success simultaneously,
except on a set of probability $o(h^{-20})$.  This includes the label
draws at a final success breaching a cap.  The bounds are uniform in
depth, cutoff, and the deterministic choices made at the states.
\end{proposition}

\begin{proof}
First bound the posterior normalizing integral from below at an
admissible past.  The phase triangle inequality gives
$A_j(v+u)\le2A_j(v)+2A_j(u)$, and the observation count gives
$2\sum_jA_j\preceq h^{130}A$ for sufficiently large $d$.
Thus
\begin{equation}\label{eq:thermal-posterior-denominator}
 \int_p\exp\left(-\sum_jA_j(z_j-y)\right)dy
 \ge c\exp(-2s_d-F(\widetilde A)).
\end{equation}
Indeed in one direction from $y_0$ a segment of length $|p|/2$ remains
inside $\overline p$.  By symmetry and the spatial variance estimate,
the density $K_{\widetilde A}(y-y_0)$ assigns this segment at least a
numerical positive mass if $C_s$ is large enough: the circular
second-moment bound puts a fixed positive mass within distance $|p|/2$
of zero, and symmetry gives half of that mass in either direction.
On this segment the
integrand is at least
$e^{-2s_d}\exp(-h^{130}A(y-y_0))$, which is at least
$e^{-2s_d}\exp(-\widetilde A(y-y_0))$.

Fix one fresh clone and integrate over its value $z'$ in the
consistency set $\{z':A(z'-Z)\le4s_d\}$.  Also use the variable
$\theta^+=V-Y$, whose independent conditional density is $K_{A^+}$.
At fixed $z',\theta^+$ expand the $2T$-moment of
\[
 \left|\sum_x\nu(x)v_x(z')\chi_{gx}(Y+\theta^+)\right|
\]
in the posterior integral.  Its weights include the extra clone
factor $e^{-A(z'-y)}$ and the spatial majorant used above.  The centered
summed form is
\[
 P=\sum_jA_j+A+S_p,
 \qquad A+S_p\preceq P\preceq\widetilde A.
\]
The upper comparison follows from the observation count and large $d$.
By \Cref{eq:thermal-envelope}, the $y$ integral before the clone
normalization and the posterior division is at most
\[
 Ce^{-F(A+S_p)}e^{-E_T(\nu)}
 \le Ce^{-F(A+S_p)}e^{-t/4}.
\]
Here the phase factors involving $\theta^+$ have modulus one, and
$|v_x(z')|\le1$, so the same coefficientwise moment bound applies.
The consistency set has length at most $e^{4s_d-F(A)}$, since
$e^{-A}\ge e^{-4s_d}$ there.  This cancels the clone normalization
$e^{F(A)}$ up to $e^{4s_d}$.  The density in $\theta^+$ integrates
to one.  Dividing by \Cref{eq:thermal-posterior-denominator}, applying
the moment tail bound at $\beta/4$, and summing over the clones gives
\[
 C N_h\exp\left(
 6s_d+F(\widetilde A)-F(A+S_p)
       +2T\log(4/\beta)-t/4\right)\le e^{-t/8}.
\]
Indeed \Cref{eq:source-7}, $t\ge\varepsilon k_b$, and
$2T\log(4/\beta)=o(t)$ absorb every positive term.  The reference
success coin is independent of $Y$ given the past and so has not
altered its posterior.

For the true-law assertion, condition on $Y$ and its entire discrete
parameter path.  At a regular success the conditional mean in
\Cref{eq:thermal-success-set} has modulus at least $\beta/2$.  The
empirical mean of the $N_h$ clone polynomials evaluated at $Y$
approximates it to error $\beta/16$ except with probability
$C/(\beta^2N_h)$, by the second-moment bound for independent bounded
complex variables.  Replacing $Y$ by $V$, and removing the common
phase $\chi_{gx_0}(V-Y)$, changes each of these polynomials by at most
\[
 \sum_x\nu(x)|1-\chi_{g(x-x_0)}(V-Y)|.
\]
Its conditional expectation is at most
$\sqrt{2\int BK_{A^+}}\le C\sqrt{t/\lambda}=O(h^{-100})$
by \Cref{eq:source-9}.  Markov's inequality makes this error at most
$\beta/16$ except with probability
$C\sqrt{t/\lambda}/\beta=o(h^{-30})$, since
$\beta\ge ch^{-32}$.  Outside these errors the clone average after
replacement has modulus at least $3\beta/8$, so at least one clone
satisfies the large-polynomial condition in \Cref{eq:source-10}.
The empirical-mean error probability is $O(h^{-56})$, also
$o(h^{-30})$.

Finally sum the own-error energies of every observation on the capped
true path, each measured in its own form.  Conditional independence
and the input-only selection rule let us apply the energy estimate in
\Cref{eq:source-7} at each draw, including final draws.  There are at
most $O(h^{121})$ observations, each of conditional expected energy
$O(h(dh^{-400}+1))$.  Their total expected energy is therefore
\[
 O(dh^{-278}+h^{122})=o(s_dh^{-30}).
\]
With exception probability $o(h^{-30})$, their total energy is at
most $s_d$.  This single event makes every past admissible by taking
$y_0=Y$.  It also gives every consistency condition, since the old
main and its clone both have type $A$ and
\[
 A(Z'-Z)\le2A(Z'-Y)+2A(Z-Y)\le4s_d.
\]
There are only $O(h)$ regular successes.  Unioning their conditional
large-polynomial failure bounds and adding the total-energy exception
proves the simultaneous $o(h^{-20})$ assertion.  The bound is over
successes rather than over all attempts; past admissibility at all
attempts followed from the single total-energy event.
\end{proof}

The numerical constants $C_1,C_2$ and the permitted small
$\varepsilon$ have now been fixed independently of the later caps.
The next section chooses those caps, then fixes $D_0,P_0,A_0$, and
finally increases the threshold for $d$. The evidence clones enter only
the likelihood estimate for regular charges. The main labels also enter
the conditional-information bounds for the spatial budgets. The realized
values of both kinds of labels do not control the spatial splitting procedure.
